\subsection{Integrable subbundles}

In this section, X denotes a manifold of class \(C^{r}\) and F a vector subbundle of class \(C^{s}\) of T(X), with r, s in \(N_{K}\) , \(s \leqslant r - 1\) .

9.3.1. Let V be a manifold of class \(C^{k}\) ( \(k \in N_{K}, k \leqslant r\) ) and let f be a morphism of class \(C^{k}\) from V into X. We say that f is an integral of F if T(f) maps T(V) into F.

A submanifold Z of class \(C^{k}\) of X is called an integral submanifold of F if the injection \(Z \rightarrow X\) is an integral of F in the above sense, i.e. if one has \(\mathrm{T}_{x}(\mathbf{Z}) \subset \mathrm{F}_{x}\) for all \(x \in Z\) .

Let \(x \in X\) . If \(Z_1\) and \(Z_2\) are two integral submanifolds of \(F\) containing \(x\) and if \(T_x(Z_1) = F_x\) , there exists a neighborhood \(U\) of \(x\) such that \(U \cap Z_2 \subset Z_1\) . If moreover \(T_x(Z_2) = F_x\) , the germs of \(Z_1\) and \(Z_2\) at \(x\) coincide.

9.3.2. We say that F is integrable if there exists a foliation Y of X of class \(C^{s}\) such that \(F = T(X, Y)\) (9.2.8). Such a foliation is then unique; it is called the integral foliation of F. For a morphism f: \(V \to X\) of class \(C^{s}\) to be an integral of F, it is necessary and sufficient that f be a morphism from V into Y.

Example. --- If F has rank 1 at every point, F is integrable, and defines on X a pure foliation of dimension 1. Suppose moreover that K = R, that X is separated, and that F admits as a frame a vector field \(\xi\) nowhere zero. Then, for every \(x \in X\) , the maximal connected leaf containing \(x\) is the image of the maximal integral arc of \(\xi\) with origin \(x\) (9.1.3).

9.3.3 ("Integrability criteria"). The following conditions are equivalent:

\begin{enumerate}
\def\labelenumi{(\roman{enumi})}
\item
  The vector subbundle F of T(X) is integrable.
\item
  For every \(x \in X\) , there exists an integral submanifold \(Z\) of class \(C^s\) of \(F\) such that \(x \in Z\) and \(T_x(Z) = F_x\) .
\item
  For every open set U of X and vector fields \(\xi\) and \(\eta\) belonging to \(\mathcal{S}_{\mathbf{F}}^{s}(\mathbf{U})\) (7.4.1), one has \([\xi, \eta] \in \mathcal{S}_{\mathbf{F}}^{s-1}(\mathbf{U})\) .
\end{enumerate}

Let \((\xi_{i})_{i\in I}\) be a family of sections of F of class \(C^{s}\) such that, for every \(x\in X\) , the set of the \(\xi_{i}(x)\) is a total subset of the Banach space \(F_{x}\) (EVT, I, § 2, no. 1). The conditions (i), (ii), and (iii) are then equivalent to:

\begin{enumerate}
\def\labelenumi{(\roman{enumi})}
\setcounter{enumi}{3}
\tightlist
\item
  for every pair \((i,j)\) of elements of I and every \(x\in X\) , one has \([\xi_i,\xi_j](x)\in F_x\) .
\end{enumerate}

When I is finite and the family \((\xi_{i})\) is a frame of F, the preceding conditions are still equivalent to:

\begin{enumerate}
\def\labelenumi{(\alph{enumi})}
\setcounter{enumi}{21}
\tightlist
\item
  there exists a family \((c_{ij}^{k})_{(i,j,k)\in I\times I\times I}\) of functions on \(X\) with values in \(K\) such that \([\xi_i,\xi_j] = \sum_{k}c_{ij}^k\xi_k\) for all \(i,j\) in \(I\) .
\end{enumerate}

(If this is the case, the functions \(c_{i,j}^{k}\) are of class \(\mathbf{C}^{s-1}\) .)

9.3.4. Let \(s' \in \mathbf{N}_{\mathbb{K}}\) , with \(s' \leqslant s\) , and let \(\mathbf{F}'\) be the vector bundle of class \(\mathbf{C}^{s}\) deduced from \(\mathbf{F}\) by weakening of structure (8.7.1). For \(\mathbf{F}\) to be integrable, it is necessary and sufficient that \(\mathbf{F}'\) be integrable. In this case, the integral foliation of \(\mathbf{F}'\) is deduced from that of \(\mathbf{F}\) by weakening of structure.

9.3.5. Let E be a Banach space and \(p\) be an integer \(\geqslant 0\) . For every \(x \in X\) , denote by

\(N(p, E)_{x}\) the vector subspace of \(\text{Alt}^{p}(T_{x}(X); E)\) consisting of the elements \(u\) such that \(u(v_{1}, \ldots, v_{p}) = 0\) for \(v_{1}, \ldots, v_{p}\) in \(F_{x}\) . The spaces \(N(p, E)_{x}\) , for \(x \in X\) , are the fibers of a vector subbundle of class \(C^{s}\) of \(\text{Alt}^{p}(T(X); E)\) . Denote it by \(N(p, E)\) . If \(F\) is integrable, we have:

\begin{enumerate}
\def\labelenumi{(\roman{enumi})}
\setcounter{enumi}{5}
\tightlist
\item
  for every open set U of X and every form \(\omega \in \mathcal{S}_{\mathrm{N}(p,\mathbb{E})}^s (\mathrm{U})\) , one has \(d\omega \in \mathcal{S}_{\mathrm{N}(p + 1,\mathbb{E})}^{s - 1}(\mathrm{U})\) .
\end{enumerate}

Conversely, if (vi) holds for \(p=1\) and for every Banach space E, the bundle F is integrable; when K = R or C, it even suffices to verify (vi) for \(p=1\) and E = K.

Suppose that the dual of T(X)/F admits a frame \((\omega_{1},\ldots,\omega_{n})\) . The integrability of F is then equivalent to the following condition:

\begin{enumerate}
\def\labelenumi{(\roman{enumi})}
\setcounter{enumi}{6}
\tightlist
\item
  \[
  d \omega_ {i} \wedge \omega_ {1} \wedge \dots \wedge \omega_ {n} = 0 \quad \text {for} \quad 1 \leqslant i \leqslant n.
  \]
\end{enumerate}

If moreover there exists a vector subbundle G of class \(C^{s}\) of T(X) such that T(X) is the direct sum of F and G, condition (vii) is equivalent to:

\begin{enumerate}
\def\labelenumi{(\roman{enumi})}
\setcounter{enumi}{7}
\tightlist
\item
  there exist differential forms \(\alpha_{i}^{j}\) ( \(i, j\) in I) of degree 1 on X, of class \(\mathbf{C}^{s-1}\) , such that \(d\omega_{i} = \sum_{j} \alpha_{i}^{j} \wedge \omega_{j}\) for all \(i \in \mathbf{I}\) .
\end{enumerate}

9.3.6. Let L be a Banach space and let \(\omega\) be a differential form of degree 1 on X with values in L, of class C \(^{s}\) . We make the following two hypotheses:

\begin{enumerate}
\def\labelenumi{\alph{enumi})}
\item
  for all \(x \in \mathbf{X}\) , \(\omega_x\) is a surjective homomorphism of \(\mathrm{T}_x(\mathbf{X})\) to \(\mathbf{L}\) , and the kernel of \(\omega_x\) is \(\mathrm{F}_x\) ;
\item
  there exists a vector subbundle G of T(X) such that T(X) is the direct sum of F and G.
\end{enumerate}

The integrability of F is then equivalent to:

\begin{enumerate}
\def\labelenumi{(\roman{enumi})}
\setcounter{enumi}{8}
\tightlist
\item
  there exists a differential form \(\alpha\) of degree 1 on X with values in \(\operatorname{End}(\mathbf{L})\) such that \(d\omega = \alpha \wedge \omega\) , the exterior product being defined by the canonical pairing \(\operatorname{End}(\mathbf{L}) \times \mathbf{L} \to \mathbf{L}\) (7.8.2 and 8.3.2).
\end{enumerate}

9.3.7 ("Equations in total differentials"). Suppose that X is the product of two manifolds A and B of class \(C^{r}\) ; one denotes \(p_{1}: X \to A\) and \(p_{2}: X \to B\) the two projections. Let f be a morphism of class \(C^{s}\) of \(p_{1}^{*}T(A)\) to \(p_{2}^{*}T(B)\) . The graphs of the maps \(f_{(a,b)}: T_{a}(A) \to T_{b}(B)\) are the fibers of a vector subbundle of class \(C^{s}\) of \(T(X)\) ; denote it by \(F^{f}\) .

Let A' be an open subset of A, and let \(\varphi: \mathbf{A}' \to \mathbf{B}\) be a morphism of class \(\mathbf{C}^k\) ( \(k \in \mathbb{N}_{\mathbb{K}}\) , \(k \leqslant r\) ). One says that \(\varphi\) is an integral of \(f\) if, for all \(a \in \mathbf{A}'\) , one has \(\mathrm{T}_a(\varphi) = f_{a,\varphi(a)}\) ; this is equivalent to saying that the map \(a \mapsto (a, \varphi(a))\) of \(\mathbf{A}'\) into X is an integral of \(\mathbf{F}^f\) (9.3.1). Such a map is of class \(\mathbf{C}^{s+1}\) . If \(\varphi_1\) and \(\varphi_2\) are two integrals of \(f\) and take the same value at a point \(a \in \mathbf{A}\) , they coincide in a neighborhood of \(a\) .

More generally, let Z be a manifold of class \(C^{k}\) , let A' be an open subset of A, let \(a \in A'\) and let \(\Phi_{1}, \Phi_{2}\) be morphisms of class \(C^{k}\) from \(Z \times A'\) into B.

\(\Phi_{1}\) and \(\Phi_{2}\) coincide on \(Z \times \{a\}\) and that, for all \(z \in Z\) , the morphisms

\[
a \mapsto \Phi_ {1} (z, a) \quad \text {and} \quad a \mapsto \Phi_ {2} (z, a)
\]

are integrals of \(f\) . Then \(\Phi_{1}\) and \(\Phi_{2}\) coincide in a neighborhood of \(Z \times \{a\}\) .

Suppose that \(F^{f}\) is integrable. Let Z be a manifold of class \(C^{k}\) ( \(k \in \mathbf{N}_{\mathbb{K}}, k \leqslant s\) ), \((z_{0}, a_{0})\) a point of \(Z \times A\) , and \(\rho\) be a morphism of class \(C^{k}\) from Z into B. There exists an open neighborhood \(Z' \times A'\) of \((z_{0}, a_{0})\) in \(Z \times A\) and a morphism \(\Phi: Z' \times A' \to B\) of class \(C^{k}\) such that, for all \(z \in Z'\) , the map \(a \mapsto \Phi(z, a)\) from \(A'\) into B admits f as its tangent map and takes the value \(\rho(z)\) at the point \(a_{0}\) .

9.3.8. We retain the hypotheses and notation of 9.3.7 and suppose that A (resp. B) is an open submanifold of a Banach space E (resp. M). The map f then identifies with a morphism of class C \(^{s}\) from X = A × B into the Banach space \(\mathcal{L}(E; M)\) . Let D \(_{1}\) f (resp. D \(_{2}\) f) denote the first (resp. second) partial derivative of f (1.6.2); it is a morphism of class C \(^{s-1}\) from X into \(\mathcal{L}(E; \mathcal{L}(E; M))\) (resp. into \(\mathcal{L}(M; \mathcal{L}(E; M))\) ), a space identified in the obvious way with \(\mathcal{L}_{2}(E; M)\) (respectively with \(\mathcal{L}(M, E; M)\) ). For F to be integrable, it is necessary and sufficient that, for every x ∈ X, the bilinear map

\[
\Delta_ {x}: (h _ {1}, h _ {2}) \mapsto \mathrm{D} _ {1} f (x) (h _ {1}, h _ {2}) + \mathrm{D} _ {2} f (x) (f (x) h _ {1}, h _ {2})
\]

from \(E \times E\) to M be symmetric. Under this condition, if \(\varphi\) is an integral of f defined in an open subset A' of A, the second derivative of \(\varphi\) at a point \(a\) of A' is \(\Delta(a,\varphi(a))\).

If \(E = K^n\) , and if one denotes \((x^1, \ldots, x^n)\) the coordinate functions on \(K^n\) , the map \(f\) is defined by a family \((f_1, \ldots, f_n)\) of maps from \(A \times B\) to \(M\) , and the integrability condition is written:

\[
\frac {\partial f _ {i}}{\partial x ^ {j}} + (\mathbf {D} _ {2} f _ {i}). f _ {j} = \frac {\partial f _ {j}}{\partial x ^ {i}} + (\mathbf {D} _ {2} f _ {j}). f _ {i}
\]

for all integers \(i, j\) in \([1, n]\) . A map \(\varphi\) of an open subset A' of A into B is an integral of \(f\) if and only if one has

\[
\frac {\partial \varphi}{\partial x ^ {i}} = f _ {i} (x, \varphi (x)) \quad \text {for all} \quad x \in \mathbf {A} ^ {\prime} \text { and all } i \in [ 1, n ],
\]

in other words, if one has

\[
d \varphi = \sum_ {1 \leqslant i \leqslant n} f _ {i} (x, \varphi (x)) d x ^ {i}.
\]

